\documentclass[a4paper,12pt,oneside]{article}

\usepackage[T2A]{fontenc}
\usepackage[utf8]{inputenc}
\usepackage[russian]{babel}

\usepackage{lmodern}
\usepackage{microtype}

\usepackage{amsmath,amssymb,amsthm,mathtools}

\usepackage{geometry}
\geometry{top=25mm,left=25mm,right=20mm,bottom=25mm}

\usepackage{graphicx}

\usepackage{enumitem}
\setlist[enumerate]{leftmargin=*, itemsep=0.35\baselineskip, topsep=0.35\baselineskip}
\setlist[itemize]{leftmargin=*, itemsep=0.2\baselineskip, topsep=0.25\baselineskip}

\newlist{alphlist}{enumerate}{1}
\setlist[alphlist]{label=\textbf{\alph*)}, leftmargin=*, align=left}

\newlist{numlist}{enumerate}{1}
\setlist[numlist]{label=\textbf{\arabic*)}, leftmargin=*, align=left}

\newcommand{\Task}[1]{%
  \subsection*{#1}%
  \vspace{-0.2\baselineskip}%
}

\setlength{\parindent}{1.25em}
\setlength{\parskip}{0pt}

\begin{document}

\begin{center}
    {\Large\bfseries Булевы функции и нормальные формы}\\[2mm]
\end{center}
\vspace{2mm}
\hrule
\vspace{4mm}

\Task{Задание 1}
Укажите номера формул, представленных в ДНФ.
\begin{numlist}
    \item $AB+CD$
    \item $A+B+C$
    \item $A+BC+\overline{E}$
    \item $P+Q(P+R)$
    \item $A+A$
\end{numlist}

\Task{Задание 2}
Укажите номера формул, представленных в КНФ.
\begin{numlist}
    \item $(AC+B)A$
    \item $A(B+C)$
    \item $B(AB+\overline{A}B)$
    \item $ABC$
    \item $A+B$
\end{numlist}

\Task{Задание 3}
Найдите инверсию выражения и упростите его.
\[
(\overline{A}+\overline{B}+\overline{C})(\overline{B}+\overline{C})(\overline{B}+\overline{C}+\overline{D})
\]

\Task{Задание 4}
Указать номера функций, представленных:\\ а) в СДНФ; б) в СКНФ; в) в ДНФ; г) в КНФ.
\begin{numlist}
    \item $f(A,B)=\overline{A}\,\overline{B}+A\overline{B}$
    \item $f(A,B,C)=AB+AC+BC$
    \item $f(A,B,C)=(A+B+C)(\overline{A}+B+\overline{C})(A+\overline{B}+C)$
    \item $f(A,B,C,D)=(\overline{A}+B+\overline{C}+D)(\overline{B}+C+D)(A+\overline{C}+D)$
    \item $f(A,B,C,D)=\overline{A}BC+\overline{B}CD+\overline{A}CD+\overline{A}BD$
    \item $f(A,B,C)=\overline{A}\,\overline{B}C+\overline{A}B\overline{C}+\overline{A}\,\overline{B}\,\overline{C}$
    \item $f(A,B,C,D)=(A+\overline{B}+C)(A+B+C)(A+\overline{B+C})$
\end{numlist}

\Task{Задание 5}
Записать десятичные индексы макстермов.
\begin{numlist}
    \item $A+\overline{B}+C+D$
    \item $R+S+\overline{T}+\overline{K}$
    \item $P+Q+\overline{R}+S$
    \item $P+\overline{Q}$
    \item $\alpha+\overline{\beta}+\overline{\gamma}$
    \item $A+\overline{B_1}+B_2$
\end{numlist}
Записать десятичные индексы минтермов, соответствующих инверсиям указанных макстермов.

\Task{Задание 6}
Укажите номера функций, возвращающих единичное значение на наборе 12.
\begin{numlist}
    \item $f=A\overline{B}+B\overline{D}+\overline{A}C$
    \item $f=BD+AC+CD$
    \item $f=\overline{D}+\overline{A}C+\overline{B}D$
    \item $f=\overline{C}+BD+\overline{A}\,\overline{B}$
    \item $f=ABC+BD$
    \item $f=A\overline{C}+AC+BD$
\end{numlist}

\Task{Задание 7}
Запишите двоичные наборы, на которых минтермы возвращают единичное значение:
\[
1.\; AB\overline{C}\,\overline{D}E
\qquad
2.\; VX\overline{Y}Z
\]

\Task{Задание 8}
Укажите номера, где записаны минтермы.
\begin{numlist}
    \item $AB\overline{C}$
    \item $ABAC$
    \item $A+\overline{B}+C$
    \item $BC\overline{D}$
    \item $PQ\overline{R}S$
    \item $ACM$
    \item $AK\overline{K}B$
    \item $ABBC$
\end{numlist}

\Task{Задание 9}
Сколько существует минтермов шести аргументов, двоичные индексы которых начинаются с нуля.

\Task{Задание 10}
Для указанной функции записать СДНФ:
\[
f_1=(01101011),\quad f_2=(01001110),\quad f_3=(0100111001)
\]

\Task{Задание 11}
Представить функцию в СДНФ при помощи таблицы истинности:
\[
f(A,B,C)=\overline{A}+B\overline{C}
\]

\Task{Задание 12}
Получить СДНФ для функции. Представить функцию в виде вектора минтермов:
\[
f=(x\oplus y)\to y\cdot x
\]

\Task{Задание 13}
Разложить функцию по всем переменным, чтобы в результате получить СДНФ:
\[
f(A,B,C,D)=AB+\overline{C},\qquad f(A,B,C,D)=A\overline{C}
\]

\Task{Задание 14}
Построить СКНФ функций из предыдущего задания по известным СДНФ.

\Task{Задание 15}
Алгоритм распознавания фиктивной переменной по таблице истинности. Перечислить все существенные и фиктивные переменные у функций:
\[
f_1=(0101111101011111),\qquad f_2=(1011010110110101)
\]
Удалить фиктивные переменные.

\end{document}
